\subsubsection{Action character and provenance of its coefficients}
\label{sec:revision-planck}

The reduced Planck constant, \(\hbar\), is the physical object
motivating the realization of the angular action section. Its position in
the electronic dependency must be specified at two levels. The determinantal
norm \(\Sigma_H\) determines the decade of the scale. The character
\(H_5\) and regional return then determine the coordinates of the
action sections. This separation presents the result
required for the electron without anticipating the complete developments of
other sectors of dimensional mechanics.

The provenance of the fifth-order character can be made explicit through
invariants already fixed on the support. In the central chart we consider
\[
 B_c=\begin{pmatrix}7&2\\2&7\end{pmatrix},
 \qquad \lambda_+=9,\quad\lambda_-=5.
\]
The calendar has length \(12\), the step-\(3\) orbit has
length \(4\), and the relevant census quotient is
\(104976/1944=54\). The character coefficients satisfy
\begin{equation}
 \frac9{16}=\frac{\lambda_+}{4^2},\quad
 \frac59=\frac{\lambda_-}{\lambda_+},\quad
 \frac7{48}=\frac{\operatorname{tr}(B_c)/2}{4\cdot12},\quad
 \frac1{54}=\frac{1944}{104976}.
 \label{eq:revision-h5-coeficientes}
\end{equation}
The assignment of these invariants to degrees \(2,3,4,5\), with the
alternating signs of \eqref{eq:el-h5}, belongs to the definition of the
analytic character used. Making it explicit is essential: the set of
cardinalities, considered without the character and without the order of its degrees,
would admit other expressions. The constructive content lies in the
composition of the support, orientation and declared analytic rule.

\subsubsection{Regional continuation and renewal equation}

The selector \(169\) determines the degree-six impulse. Subsequent
prolongation is organized by the transfer \(D_A\) and a unit
normalization of the determinantal line. These rules admit a formulation
in formal series before substituting \(z=\alpha\), with \(D_A\) already
constructed. Write
\[
 \mathscr C(z)=169z^6+R(z),\qquad R(z)\in z^7\mathbb R[[z]].
\]
Concatenating a prolongation increases the degree by one and
multiplies its weight by \(D_A\). The first strict prolongation has
weight \(D_A\). Its renewal equation is therefore
\begin{equation}
 R(z)=D_Az^7+D_AzR(z).
 \label{eq:revision-renovacion}
\end{equation}

\begin{proposition}[Formal uniqueness of the normalized continuation]
Equation \eqref{eq:revision-renovacion} determines a unique series, and
\[
 \mathscr C(z)=169z^6+\frac{D_Az^7}{1-D_Az}
             =169z^6+\sum_{n=7}^{\infty}D_A^{n-6}z^n.
\]
The realization is analytic on \(|D_Az|<1\).
\end{proposition}
\begin{proof}
The element \(1-D_Az\) has constant term one and is invertible
in the formal power series ring. Solving the equation produces the displayed
series. Equivalently, the degree-seven coefficient is \(D_A\)
and each subsequent coefficient is \(D_A\) times its predecessor. The
geometric series converges absolutely in the indicated disk.
\end{proof}

Normalization of the first prolongation is an indispensable operational
datum. If only the degree-six impulse and
concatenation ratio were preserved, the series
\[
 169z^6+\lambda\frac{D_Az^7}{1-D_Az}
\]
would still share both data for any \(\lambda\).
The unit rule fixes \(\lambda=1\) before evaluation.
This locates each rule's contribution to uniqueness of the
readout, and the rational expression preserves the complete continuation.

\subsubsection{Reduced Planck constant and return of the action section}

In the dimensional chart of \eqref{eq:el-secciones-accion}, the
pre-return section is the model's construction intended for comparison
with the reduced Planck constant:
\[
 \hbar_{\mathrm{pre}}=H_5\,10^{-34}\mathcal U_S.
\]
The subsequent section also preserves the regional compensation
\(90\mathscr C(\alpha)/\pi\). Both belong to the same action
line, and their quotient \(R_{\mathrm{act}}\) measures the multiplicative
effect of return. Passage to \(h\) through \(2\pi\) uses the relation
between angular parametrization and a complete winding.

This construction effectively enters
\(R_{\mathrm{act}}^{80-54\alpha+6\Delta_4}\). The electron preserves
its central coordinate and trajectory; the chronological scale does not
replace that structure with an undifferentiated quantum. Nor is a mass
already expressed in MeV multiplied again by \(10^{-34}\):
the decade belongs to the action section, cancels in its dimensionless
quotient and reappears only where the corresponding dimensional
realization acts.

The numerical comparison of \(H_5\), the returned section and
\(h/(2\pi)\) is presented in Section~\ref{sec:revision-planck-contraste}.
This distinguishes the name of the physical object, the equation the
model assigns to it and the precision with which that equation reproduces it.
